\Needspace{7\baselineskip}
\begin{proof}[Proof of Proposition~\ref{prop:rewrite-uncapped-complements-bounds}]
\phantomsection
\label{proof:rewrite-uncapped-complements-bounds}
Equations~\eqref{eq:rewrite-uncapped-complements-z-bound}
and~\eqref{eq:rewrite-uncapped-complements-y-bound} follow because
$\sigma/(\sigma-1)<0$. Define $k=K/(AL)$. The resource constraint
\eqref{eq:rewrite-eq-kdot}, nonnegative consumption and compute expenditure,
and $L=N$ imply
\begin{equation}
 \dot k\leq
 \omega_L^{\sigma(1-\alpha)/(\sigma-1)}k^\alpha
 -(\delta+n+\gamma)k.
 \label{eq:rewrite-uncapped-complements-k-bound}
\end{equation}
Since $\delta+n+\gamma>0$ and $0<\alpha<1$, the right-hand side is
negative above a finite positive level. In particular,
\[
 k_t\leq
 \max\left\{k_0,
 \left[\frac{\omega_L^{\sigma(1-\alpha)/(\sigma-1)}}
 {\delta+n+\gamma}\right]^{1/(1-\alpha)}\right\}.
\]
Equation~\eqref{eq:rewrite-uncapped-complements-y-bound} then bounds
$Y/(AL)$ uniformly. Since $A_t=A_0e^{\gamma t}$ and $L=N$,
$Y/N$ is bounded above by a constant times $e^{\gamma t}$.
Taking logarithms, dividing by $t$, and taking the upper limit gives
\eqref{eq:rewrite-uncapped-complements-growth-bound}.

The same bounds keep $K$ and $Y$ finite on every bounded time interval.
For any such horizon $T$, integrating the resource constraint gives
\[
 \int_0^T M_t\dd t\leq K_0+\int_0^T Y_t\dd t<\infty.
\]
The uncapped research law and H\"older's inequality imply
\begin{align}
 B_T^{1-\eta}
 &=B_0^{1-\eta}+(1-\eta)\chi\int_0^T M_t^\eta\dd t
 \nonumber\\
 &\leq B_0^{1-\eta}
 +(1-\eta)\chi T^{1-\eta}
 \left(\int_0^T M_t\dd t\right)^\eta<\infty.
 \label{eq:rewrite-uncapped-complements-b-bound}
\end{align}
Thus $B$ cannot diverge at a finite date either. This argument bounds
stocks and output; it is not a regularity theorem for all instantaneous
control choices or an equilibrium-existence proof.

There is also no finite-horizon unbounded-payoff argument of the kind used
for substitution in Proposition~\ref{prop:rewrite-research-scale}.
At given positive continuous paths of $K,A,L$ on $[0,T]$, optimized operating
profit satisfies
\[
 0\leq\pi_t(B)\leq p_XX\leq(1-\alpha)Y
 \leq(1-\alpha)\omega_L^{\sigma(1-\alpha)/(\sigma-1)}
 K_t^\alpha(A_tL_t)^{1-\alpha}.
\]
This upper bound is independent of the developer's choice of $B$.
For a given locally integrable $r$, the discount factor
$D_t=\exp(-\int_0^t r_s\dd s)$ has a positive minimum on $[0,T]$.
Discounted operating revenue is bounded, whereas the cost of research
tends to infinity with $\int_0^T M_t\dd t$. The developer's truncated
supremum is therefore finite, and net payoff tends to minus infinity as
total research expenditure diverges. This conclusion holds for every
$0<\eta<1$, but does not prove that an infinite-horizon optimum exists.
\end{proof}

\Needspace{7\baselineskip}
\begin{proof}[Proof of Proposition~\ref{cond:rewrite-uncapped-complements-limits}]
\phantomsection
\label{proof:rewrite-uncapped-complements-limits}
Within this proof, write $x=X/(AL)$, $k=K/(AL)$, $c=C/(AL)$,
$y=Y/(AL)$, $z=Y/K$, and $\varphi=(\sigma-1)/\sigma<0$.

\textit{Step 1. Static optimality identifies the limiting technology.}
Equation~\eqref{eq:rewrite-uncapped-complements-k-bound} bounds $k$ above.
The assumed finite limit of $g_C$ and the Euler and capital-return equations
\eqref{eq:rewrite-euler} and~\eqref{eq:rewrite-capital-return} imply finite
limits of $r$ and $z=(r+\delta)/\alpha$.
The static monopoly condition~\eqref{eq:rewrite-monopoly-foc} requires
$1-e_X>0$. Since $e_X=(1-s_X)/\sigma+\alpha s_X$, this gives
$s_X>(1-\sigma)/(1-\alpha\sigma)>0$.
Production and pricing, as in equations
\eqref{eq:rewrite-ai-output-limit}--\eqref{eq:rewrite-output-capital-limit},
give the exact identity
\[
 z^{\alpha/(1-\alpha)}
 =B(1-\alpha)\omega_X^{\sigma/(\sigma-1)}
 (1-e_X)s_X^{-1/(\sigma-1)}.
\]
As $B\to\infty$ and $z$ remains bounded, this identity forces $e_X\to1$.
Hence $s_X\to s_\infty$ and $x\to x_\infty$, where the CES formula gives
\begin{align}
 s_\infty&=\frac{1-\sigma}{1-\alpha\sigma},&
 x_\infty&=
 \left[\frac{\omega_Ls_\infty}
 {\omega_X(1-s_\infty)}\right]^{1/\varphi}>0.
 \label{eq:rewrite-uncapped-complements-x-limit}
\end{align}
Both are finite for each fixed $\sigma<1$.
Let $f(x)=[\omega_L+\omega_Xx^\varphi]^{(1-\alpha)/\varphi}$.
Then $y=k^\alpha f(x)$ and $z=k^{\alpha-1}f(x)$.
Because $k$ is bounded above and $f(x)\to f(x_\infty)>0$,
the limit of $z$ is strictly positive. Hence
$k\to k_\infty>0$ and $y\to y_\infty>0$.
Also $U/(AL)=x/B\to0$, so $U/Y\to0$.
No convergence of capital or consumption ratios has been assumed.

\textit{Step 2. Both transversality conditions pin down consumption growth.}
If the limit of $g_C$ exceeded $n+\gamma$, then $c\to\infty$.
The normalized resource constraint would give
$\dot k\leq y-c-(\delta+n+\gamma)k$, contradicting positive finite capital.
If that limit were smaller than $n+\gamma$, then $c\to0$ and $C/Y\to0$.

To exclude the latter case, use the research first-order condition
\eqref{eq:rewrite-research-foc} and uncapped costate
equation~\eqref{eq:rewrite-costate} to obtain
\begin{equation}
 \frac{\mathrm d(qB)}{\mathrm dt}
 =rqB-U+\frac{1-\eta}{\eta}M.
 \label{eq:rewrite-complements-shadow-stock}
\end{equation}
For $D_t=\exp(-\int_0^t r_s\dd s)$, combining this equation with the
resource constraint~\eqref{eq:rewrite-resource} cancels research expenditure:
\begin{equation}
 \frac{\mathrm d}{\mathrm dt}
 \left[D_t\left(K_t+\frac{\eta}{1-\eta}q_tB_t\right)\right]
 =D_t\left[(1-\alpha)Y_t-C_t-\frac{U_t}{1-\eta}\right].
 \label{eq:rewrite-complements-weighted-tvc}
\end{equation}
The expression in parentheses on the left is an auxiliary positive
combination, not a definition of the developer's market value.
The transversality conditions~\eqref{eq:rewrite-household-tvc}
and~\eqref{eq:rewrite-developer-tvc} imply that the bracketed expression tends
to zero: Euler makes $D_t$ proportional to $e^{-\rho t}N_t/C_t$.
If $C/Y\to0$ and $U/Y\to0$, its derivative would instead be strictly
positive at all sufficiently late dates. A positive increasing expression
cannot tend to zero. Thus
\begin{equation}
 g_C\to n+\gamma,\qquad r\to\rho+\gamma,\qquad
 \frac KY\to\frac{\alpha}{\rho+\gamma+\delta}.
 \label{eq:rewrite-complements-euler-limit}
\end{equation}
The cancellation in~\eqref{eq:rewrite-complements-weighted-tvc} works for
every $0<\eta<1$; it requires neither $\eta\leq1/2$ nor $\eta<\alpha$.

\textit{Step 3. Optimal research expenditure vanishes relative to output.}
Integrate~\eqref{eq:rewrite-complements-shadow-stock} using the developer's
transversality condition~\eqref{eq:rewrite-developer-tvc}.
Since $qB>0$ and $M\geq0$,
\begin{equation}
 \int_t^\infty\frac{D_s}{D_t}M_s\dd s
 \leq\frac{\eta}{1-\eta}
 \int_t^\infty\frac{D_s}{D_t}U_s\dd s.
 \label{eq:rewrite-complements-research-pv-bound}
\end{equation}
Here the integrals exist: $r\to\rho+\gamma>n+\gamma$ and $U/(AL)\to0$
bound the discounted inference integral; integration over finite horizons
and the developer TVC then bound the nonnegative research integral.
More explicitly, put $d=(\rho-n)/2>0$. For sufficiently late $t$,
$r_s\geq n+\gamma+d$ for all $s\geq t$, so
\[
 \frac1{A_tL_t}\int_t^\infty\frac{D_s}{D_t}U_s\dd s
 \leq\frac1d\sup_{s\geq t}\frac{U_s}{A_sL_s}\longrightarrow0.
\]

This integral limit also gives a pointwise limit for $M/(AL)$.
Differentiating the research condition and using the costate equation gives
\begin{equation}
 (1-\eta)g_M
 =r-\frac{\eta\chi UB^{\eta-1}}{M^{1-\eta}}
 \leq r.
 \label{eq:rewrite-complements-research-growth-bound}
\end{equation}
Choose a positive constant $R$ above the eventual values of $r$, and write
$H=R/(1-\eta)$. Then $M_s\geq M_t e^{-H(t-s)}$ for $s\in[t-1,t]$
at sufficiently late dates. Thus
\[
 \frac1{A_{t-1}L_{t-1}}
 \int_{t-1}^t\frac{D_s}{D_{t-1}}M_s\dd s
 \geq e^{n+\gamma-R-H}\frac{M_t}{A_tL_t}.
\]
The left side tends to zero by~\eqref{eq:rewrite-complements-research-pv-bound}.
Consequently $M/(AL)\to0$ and $M/Y\to0$.

\textit{Step 4. Consumption and capital ratios converge to positive constants.}
The resource constraint now reads
\[
 \dot k=a_t-c,\qquad
 a_t:=y-\frac{U+M}{AL}-(\delta+n+\gamma)k
 \longrightarrow a_\infty:=y_\infty-(\delta+n+\gamma)k_\infty.
\]
Also $\dot c/c=g_C-n-\gamma\to0$.
Integrating over $[t,t+1]$ shows that
$\int_t^{t+1}c_s\dd s\to a_\infty$, since $k$ converges.
The vanishing logarithmic derivative implies, uniformly on this interval,
$c_s/c_t\to1$. Therefore $c_t\to a_\infty$.
Using~\eqref{eq:rewrite-complements-euler-limit},
\[
 \frac CY\to\frac{a_\infty}{y_\infty}
 =1-\frac{\alpha(\delta+n+\gamma)}{\rho+\gamma+\delta}>0.
\]
Positivity follows from $\rho>n$, $0<\alpha<1$, and $\gamma>0$.
In particular, $c$ converges to a positive constant and $\dot k\to0$.
This proves~\eqref{eq:rewrite-uncapped-complements-allocation-ratios}.

\textit{Step 5. Output and wage growth converge as well.}
A convergent level need not have a convergent derivative, so this last
step verifies the growth limits rather than inferring them from levels.
The same bound on $g_M$ used in Step 3 gives
\[
 B_t^{1-\eta}
 \geq(1-\eta)\chi M_t^\eta\int_0^1e^{-\eta H v}\dd v.
\]
The research law therefore bounds $g_B=\chi M^\eta/B^{1-\eta}$ above.
Write the static marginal-revenue condition as
$k^\alpha v(x)=1/B$, where
$v(x)=(1-\alpha)s(x)f(x)[1-e_X(x)]/x$.
At $x_\infty$, $v$ has a simple zero:
$v(x_\infty)=0$ and $v'(x_\infty)<0$, because $e_X$ increases through one.
Differentiation gives
\[
 \dot x=-\frac{v(x)}{v'(x)}
 \left(g_B+\alpha\frac{\dot k}{k}\right)\longrightarrow0.
\]
Hence $\dot y\to0$ and $\dot s_X\to0$.
Since $y\to y_\infty>0$ and $1-s_\infty>0$,
$g_Y\to n+\gamma$ and, by the wage equation~\eqref{eq:rewrite-wage},
$g_w\to\gamma$.
Substitution of $s_\infty$ into~\eqref{eq:rewrite-wage}
and~\eqref{eq:rewrite-ai-price} gives the stated limiting income shares.
This is a characterization of any equilibrium satisfying the two remaining
asymptotic hypotheses, not an equilibrium-existence theorem.
\end{proof}
